\section{Selected Controlled Results}
\label{app:controlled-results}

This appendix expands the behavioral validation in Section~\ref{sec:controlled} without repeating its main tests. We report five positive extensions: a matched-magnitude example, the complete response atlas, all evidence-training trajectories, fragility checks across interventions and reveal geometries, and contradiction checks across regimes, architectures, seeds, and counterfactual maps.

\subsection{Matched-magnitude response roles}
\label{app:controlled-overview}

Figure~\ref{fig:controlled-overview} gives a concrete example of the scalar ambiguity discussed in the main text. In an object-shape model, the true shape factor and an endpoint-null floor-color factor produce nearly equal ordinary response magnitude, but their response roles differ sharply. The same distinction matters across the controlled benchmark: ordinary magnitude can rank an endpoint-null factor above a supported factor, whereas endpoint-oriented routing removes these false-null reversals.


\begin{figure*}[t]
    \centering
    \begin{subfigure}[t]{0.46\textwidth}
        \centering
        \includegraphics[
            height=3.75cm,
            width=\linewidth,
            keepaspectratio
        ]{figures/controlled/fig1a_response_components_bar.pdf}
        \caption{Nearly equal magnitude, different response roles.}
    \end{subfigure}
    \hspace{0.04\textwidth}
    \begin{subfigure}[t]{0.30\textwidth}
        \centering
        \includegraphics[
            height=3.35cm,
            width=\linewidth,
            keepaspectratio
        ]{figures/controlled/fig1b_false_null_order.pdf}
        \caption{Orientation removes false-null reversals.}
    \end{subfigure}
    \caption{\textbf{Equal response magnitude can hide different behavior.}
    (a) The true shape factor and an endpoint-null floor-color factor have nearly identical ordinary magnitude, yet \decaf separates evidence from fragility.
    (b) Across the benchmark, ordinary magnitude ranks an endpoint-null factor above a supported factor in $7.32\%$ of comparable pairs, whereas endpoint-oriented routing reduces this rate to zero.}
    \label{fig:controlled-overview}
\end{figure*}

\subsection{Complete decomposition atlas}
\label{app:controlled-atlas}

Figure~\ref{fig:app-controlled-atlas} displays every task--factor combination. Endpoint-null factors concentrate near the fragility corner, direct color tasks concentrate in evidence, and interaction tasks carry more contradiction and mixed mass. Across 125 endpoint-null units, 96.8\% of ordinary response is fragility on average. Across 48 endpoint-supported units, the mean composition is 64.1\% evidence, 12.5\% contradiction, and 23.4\% null fragility.

\begin{figure*}[t]
    \centering
    \includegraphics[width=0.7\textwidth]{figures/appendix/controlled/figA1_decomposition_atlas.pdf}
    \caption{\textbf{Complete controlled decomposition atlas.}
    Results cover five tasks, two architectures, three seeds, and six factors. Endpoint-null factors are dominated by fragility, while task-relevant and interaction factors contain evidence and contradiction in different proportions.}
    \label{fig:app-controlled-atlas}
\end{figure*}

\subsection{Evidence across training trajectories}
\label{app:e-trajectories}

Figure~\ref{fig:app-e-heatmaps} places exact shortcut vulnerability and the \decaf evidence margin on the same eight checkpoint trajectories. Six trajectories exhibit nonconstant vulnerability, and all six show positive within-trajectory association between evidence margin and shortcut vulnerability. The two Small-ViT trajectories trained at the strongest shortcut correlation remain shortcut-dominant, and \decaf correspondingly keeps their evidence margin high.

\begin{figure*}[t]
    \centering
    \begin{subfigure}[t]{0.47\textwidth}
        \centering
        \includegraphics[
            height=3.05cm,
            width=\linewidth,
            keepaspectratio
        ]{figures/appendix/controlled/fig2a_vrev_heatmap.pdf}
        \caption{Exact shortcut vulnerability.}
    \end{subfigure}
    \hfill
    \begin{subfigure}[t]{0.47\textwidth}
        \centering
        \includegraphics[
            height=3.05cm,
            width=\linewidth,
            keepaspectratio
        ]{figures/appendix/controlled/fig2b_e_margin_heatmap.pdf}
        \caption{Evidence margin.}
    \end{subfigure}
    \caption{\textbf{Evidence follows shortcut reliance across complete training trajectories.}
    The two panels show the same eight checkpoint sequences through an external behavioral measure and the \decaf evidence margin.}
    \label{fig:app-e-heatmaps}
\end{figure*}

\subsection{Fragility: intervention separation and reveal geometry}
\label{app:f-extensions}

The main text validates $F$ against independently measured off-path prediction change. Figure~\ref{fig:app-f-extensions} provides three complementary checks. First, $F$ cleanly orders robust, neutral, and fragile interventions. Second, similar total response can correspond to different endpoint-relative roles across architectures. Third, the robust--neutral--fragile ordering transfers from the primary covariance-matched reveal to diagonal covariance, trace-matched pixel Gaussian, and three covariance-power geometries. The controlled fragility result is therefore not tied to one second-order reveal geometry, although the numerical summaries remain protocol-relative.

\begin{figure*}[t]
    \centering
    \begin{subfigure}[t]{0.27\textwidth}
        \centering
        \includegraphics[
            height=3.5cm,
            width=\linewidth,
            keepaspectratio
        ]{figures/controlled/fig3a_f_variant_separation.pdf}
        \caption{Intervention separation.}
    \end{subfigure}
    \hfill
    \begin{subfigure}[t]{0.27\textwidth}
        \centering
        \includegraphics[
            height=3.5cm,
            width=\linewidth,
            keepaspectratio
        ]{figures/controlled/fig3c_f_same_abs_different_mechanism.pdf}
        \caption{Similar magnitude, different roles.}
    \end{subfigure}
    \hfill
    \begin{subfigure}[t]{0.38\textwidth}
        \centering
        \includegraphics[
            height=3.0cm,
            width=\linewidth,
            keepaspectratio
        ]{figures/appendix/controlled/figA10_geometry_transfer_heatmap.pdf}
        \caption{Cross-geometry rank transfer.}
    \end{subfigure}
    \caption{\textbf{Additional fragility checks.}
    (a) $F$ separates robust, neutral, and fragile training.
    (b) Similar ordinary response can be endpoint evidence in one model and endpoint-null fragility in another.
    (c) The controlled ordering transfers across held-out reveal geometries.}
    \label{fig:app-f-extensions}
\end{figure*}
\begin{figure*}[!t]
    \centering

    % ----- Regime-level diagnostics -----
    \begin{subfigure}[t]{0.285\textwidth}
        \centering
        \includegraphics[
            height=3cm,
            width=\linewidth,
            keepaspectratio
        ]{figures/controlled/fig4a_c_mechanism_summary.pdf}
        \caption{Attenuation vs.\ inversion.}
    \end{subfigure}
    \hfill
    \begin{subfigure}[t]{0.285\textwidth}
        \centering
        \includegraphics[
            height=3cm,
            width=\linewidth,
            keepaspectratio
        ]{figures/controlled/fig4b_c_label_behavior.pdf}
        \caption{Independent label behavior.}
    \end{subfigure}
    \hfill
    \begin{subfigure}[t]{0.285\textwidth}
        \centering
        \includegraphics[
            height=2.7cm,
            width=\linewidth,
            keepaspectratio
        ]{figures/controlled/fig4d_c_vs_abs_corr.pdf}
        \caption{$C$ succeeds where $\Abs$ fails.}
    \end{subfigure}

    \vspace{5pt}

    % ----- Calibration and transfer -----
    \begin{subfigure}[t]{0.285\textwidth}
        \centering
        \includegraphics[
            height=3cm,
            width=\linewidth,
            keepaspectratio
        ]{figures/appendix/controlled/figA11_c_phi_calibration.pdf}
        \caption{Opposed-fraction calibration.}
    \end{subfigure}
    \hfill
    \begin{subfigure}[t]{0.285\textwidth}
        \centering
        \includegraphics[
            height=3cm,
            width=\linewidth,
            keepaspectratio
        ]{figures/appendix/controlled/figA12_c_seed_heatmap.pdf}
        \caption{Architecture and seed consistency.}
    \end{subfigure}
    \hfill
    \begin{subfigure}[t]{0.285\textwidth}
        \centering
        \includegraphics[
            height=3cm,
            width=\linewidth,
            keepaspectratio
        ]{figures/appendix/controlled/figA13_c_wallmap_transfer.pdf}
        \caption{Counterfactual-map transfer.}
    \end{subfigure}

    \caption{\textbf{Additional contradiction checks.}
    Top: contradiction separates attenuation from inversion, agrees with independently measured label behavior, and succeeds where ordinary magnitude does not.
    Bottom: the opposed fraction is calibrated and stable across architectures, seeds, and alternative wall-color counterfactual maps.}
    \label{fig:app-c-extensions}
\end{figure*}

\subsection{Contradiction: regime separation and calibration}
\label{app:c-extensions}



The main text shows that $C$ tracks independently measured label reversal. Figure~\ref{fig:app-c-extensions} expands this result. Direct preserves the effect, Gate suppresses it, and Invert reverses it; only inversion creates substantial opposed mass. The induced label behavior independently separates into preservation, collapse, and swap. Across architectures, seeds, and wall-color maps, the opposed fraction remains well calibrated. Raw sign changes can be frequent when the stage effect is numerically close to zero, but the corresponding opposed mass remains small, separating magnitude-aware contradiction from sign counting.


